\uuid{NS04}
\titre{Matrice d'une application linéaire/Noyau/Image/Changement de bases}
\auteur{hanani ; mzari}
\organisation{exo7}
\licence{CC BY-NC-SA 4.0}
\datecreate{2020-01-20}
\niveau{L1}
\module{Algèbre}
\chapitre{Matrice}
\sousChapitre{Matrice et application linéaire}
\difficulte{3}
\theme{Matrice d'une application linéaire, Noyau, Image, Changement de bases}
\exo7qcmid{lille-2.matrice-application.018}

\contenu{
\begin{question}
\qclassification{108.03}
\qtags{motcle=Matrice d'une application linéaire/Noyau/Image/Changement de bases}
Soit $f$ l'endomorphisme de $\Rr^3$ dont la matrice dans la base canonique ${\cal B}=\{e_1,e_2,e_3\}$ est :
$$A = \left(\begin{array}{rcc}
1&-1&1\\
0&1&0\\
0&-1&2\\
\end{array}\right).$$
Quelles sont les assertions vraies ?
\begin{answers}
\good{$\dim \ker (f-Id) = 2 $ et $\dim \ker (f-2Id) = 1 $}.
\bad{$\dim \ker (f-Id) = 1 $ et $\dim \ker (f-2Id) = 2 $}.
\good{Il existe une base de $\Rr^3$ dans laquelle la matrice de $f$ est : $B= \left(\begin{array}{rcc}
1&0&0\\
0&1&0\\
0&0&2\\
\end{array}\right)$}.
\good{Il existe une matrice $C$ inversible telle que : $C^{-1}AC=  \left(\begin{array}{rcc}
1&0&0\\
0&1&0\\
0&0&2\\
\end{array}\right)$}.
\end{answers}
\begin{explanations} On vérifie que $\dim \ker (f-Id) = 2 $ et $\dim \ker (f-2Id) = 1 $. Soit $\{u_1,u_2\}$ une base de $\ker (f-Id)$ et  $\{u_3\}$ une base de $\ker (f-2Id)$.
\vskip0mm
Soit $\lambda_1, \lambda_2,\lambda_3$ des réels tels que $\lambda_1u_1+\lambda_2u_2+\lambda_3u_3=0$. En considérant l'image par $f$, on obtient $\lambda_1u_1+\lambda_2u_2+2\lambda_3u_3=0$. On déduit que $\lambda_3=0$ et comme $\{u_1,u_2\}$  est libre,
$\lambda_1=\lambda_2=0$.
\vskip0mm
Par conséquent, ${\cal B}' = \{u_1,u_2,u_3\}$ est une base de $\Rr^3$. Dans cette base, la matrice de $f$ est : $B= \left(\begin{array}{rcc}
1&0&0\\
0&1&0\\
0&0&2\\
\end{array}\right)$. On prend $C$ la matrice de passage de la base  ${\cal B}$ à la base ${\cal B}'$, c.à.d la matrice de l'identité de $\Rr^3$ de la base ${\cal B}'$ à la base  ${\cal B}$.
\end{explanations}
\end{question}
}
